\documentclass[11pt,letterpaper]{article}
\usepackage{preamble}
\authorsig{KUKLA}
\begin{document}
\moduletitle{Q32 --- Discussion}{Quantifying information flow with entropy}

\begin{enumerate}[leftmargin=*,itemsep=0.6em]
\item \textbf{Entropy vs.\ correctness.} A model answers a factual question with
very low next-token entropy but is wrong. A second model answers with high entropy
and is right. What does each case tell you, and why is ``confident'' not the same
as ``calibrated''? (Connect to Q11/Q12.)

\item \textbf{The irreducible term.} Cross-entropy decomposes into
$H(p_\text{data}) + D_\text{KL}(p_\text{data}\|p_\theta)$. If a lab reports a new
model with lower perplexity on the same corpus, which term moved --- and could a
reported perplexity gain ever be an artefact of the \emph{tokenizer} rather than
the model? How would you control for that?

\item \textbf{The information-plane debate.} Shwartz-Ziv \& Tishby saw a
``compression phase''; Saxe et al.\ showed it depends on the nonlinearity and the
binning. If a classmate shows you an information-plane plot of a network, what
three questions do you ask before believing it?

\item \textbf{Attention entropy as interpretation.} A given head has near-zero
entropy on syntax tokens but near-maximal entropy on content words. Is this head
``doing more'' on syntax? What alternative explanations (induction behaviour,
positional averaging, a ``no-op''/attention-sink head) could produce the same
entropy signature, and how would you distinguish them?

\item \textbf{Estimate vs.\ intervene.} You want to know whether information about
token $i$ reaches the prediction at position $j$. One student proposes estimating
$I(h_i; \hat y_j)$ with MINE; another proposes activation patching. Under what
conditions does each answer a \emph{different} question, and which one answers the
\emph{causal} question you probably care about?
\end{enumerate}

\end{document}
